\section{相关工作}

\subsection{非侵入式词级解码与神经上下文}

Défossez 等人通过对比学习将 EEG/MEG 与预训练语音表示对齐，
在自然连续语音条件下实现了语音片段检索\cite{defossez2023decoding}。
d'Ascoli 等人进一步将词级神经活动映射到语言表示空间，
并利用跨词上下文提高单词检索性能\cite{dascoli2025word}。
Brain2Qwerty 在打字句子解码中先提取局部神经特征，
再建模字符之间的关系，并结合语言模型改善文本\cite{levy2026brain2qwerty}。
这些工作表明，目标之外的序列上下文有助于预测，
但对上下文范围与输入可用时刻的联合比较仍较有限。
SimpleB2T 进一步指出，词对齐的重叠窗口可能提供词间时间线索\cite{jayalath2026timing}。
本文不试图识别跨词性能提升的具体神经机制，而是在已知词级对齐的既定条件下，
比较局部上下文范围及其对应的输入等待，并研究有限范围内的连续预测修正。

\subsection{增量神经解码}

侵入式语言脑机接口中，随神经信号到达而更新输出已有较多研究。
Moses 等人在新的词级证据出现后结合语言模型调整此前的文本假设
\cite{moses2021neuroprosthesis}；
Willett 等人持续更新音素概率并进行语言模型搜索，实现流式文本解码
\cite{willett2023speech}；
脑到语音系统则将更新间隔进一步缩短\cite{littlejohn2025streaming}。
文本搜索可以随着新证据到达改变历史假设。
本文关注目标神经表示的更新：后续神经活动直接参与较早目标的表示计算，
使同一目标在获得更多输入后产生新的候选分布。
非侵入式方面，CogReader 将 fMRI 序列分段，
把已解码内容的语义表示传递给后续片段\cite{lu2025braininspired}，
主要利用历史辅助未来预测。

\subsection{同步机器翻译与流式输出}

在输入等待与质量之间取舍，是同步机器翻译和流式语音识别中的核心问题。
wait-$k$ 策略先读入 $k$ 个源词，再交替读入与输出，使目标生成保持固定的源端词数滞后
\cite{ma2019stacl}。
重翻译方法允许修改已输出的内容，并用归一化擦除指标衡量显示稳定性
\cite{arivazhagan2020retranslation}。
本文同样限制需要等待的后续输入，并允许近期预测继续修正：
先依据当前输入给出预测，再在至多两个后续词窗内更新目标神经表示与文本，随后固定。
